\documentclass[11pt]{article}

\usepackage[margin=1in]{geometry}
\usepackage{amsmath,amssymb,amsthm}
\usepackage{graphicx}
\usepackage[hidelinks]{hyperref}

\newtheorem{theorem}{Theorem}

\newcommand{\C}{\mathbb C}
\newcommand{\Z}{\mathbb Z}
\newcommand{\cF}{\mathcal F}
\newcommand{\cC}{\mathcal C}

\title{A Non-Positive Kazhdan Projection for a Banach-Space Representation Family}
\author{arXiv automatic math run \texttt{fa\_banach\_001}}
\date{June 28, 2026}

\begin{document}
\maketitle

\section*{Status}

This packet gives a candidate counterexample to the general question in
Section 3.5 of de la Salle \cite{deLaSalle2016}: if
\(\cC_{\cF}(G)\) contains a Kazhdan projection, must that projection belong to
the closure of the nonnegative compactly supported functions of integral one?

The answer is no for arbitrary Banach-space representation families. The
example is finite-dimensional and non-isometric. It does not contradict the
source paper's positive result for central Kazhdan projections attached to
isometric representation families.

\section*{Source Question}

In Section 3.5, after defining the positive variant of a Kazhdan projection,
the source paper states:

\begin{center}
\includegraphics[width=\textwidth]{figures/open_problem_crop.png}
\end{center}

The question is whether, in general, the existence of a Kazhdan projection
forces the existence of one in the closure of nonnegative compactly supported
probability functions. For the discrete group used below, these functions are
exactly finitely supported probability measures on \(\Z\).

\section*{Definitions Used}

Let \(G\) be a locally compact group and let \(\cF\) be a set of
Banach-space representations satisfying the uniform generator-bound hypothesis
of the source. The seminorm on \(C_c(G)\) is
\[
  \|a\|_{\cF}=\sup_{(\pi,E)\in\cF}\|\pi(a)\|_{B(E)},
\]
and \(\cC_{\cF}(G)\) is the corresponding completion.

A Kazhdan projection is an idempotent \(p\in\cC_{\cF}(G)\), lying in the
closure of functions with integral one, such that \(\pi(p)\) is the projection
onto \(E^\pi\) for every \((\pi,E)\in\cF\). The positive variant asks that
\(p\) lie in the closure of the nonnegative functions with integral one.

\section*{Proof Intuition}

The ordinary Kazhdan projection only asks for approximation by signed
mass-one functions. A signed Laurent polynomial can annihilate two spectral
values on opposite sides of the unit circle at the same time. Positive
probability measures cannot do this: if \(N\) is an integer-valued random
variable, the two moments \(E(2^N)\) and \(E(2^{-N})\) cannot both be small
because \(2^N+2^{-N}\geq 2\) pointwise. This separates signed Kazhdan
projections from positive ones in the smallest possible diagonal example.

\begin{theorem}
There exist a compactly generated infinite group \(G\), a family \(\cF\) of
Banach-space representations satisfying the source hypotheses, and a Kazhdan
projection in \(\cC_{\cF}(G)\) which is not in the closure of the nonnegative
compactly supported probability functions.
\end{theorem}

\begin{proof}
Let \(G=\Z\), written multiplicatively with generator \(u\). Let
\(E=\C^3\) with the supremum norm, and define a representation
\(\pi\colon \Z\to GL(E)\) by
\[
  \pi(u)=T=\operatorname{diag}(1,2,1/2).
\]
Let \(\cF=\{(\pi,E)\}\). If \(S=\{-1,0,1\}\), then
\(\sup_{s\in S}\|\pi(s)\|=2\), so the source's generator-bound hypothesis is
satisfied.

For a finitely supported function \(a=\sum_n a_n u^n\in \C[\Z]\), write
\[
  \widehat a(\lambda)=\sum_n a_n\lambda^n.
\]
Then
\[
  \pi(a)=\operatorname{diag}\bigl(\widehat a(1),
  \widehat a(2),\widehat a(1/2)\bigr).
\]
The invariant vectors of \(\pi\) are exactly
\[
  E^\pi=\C e_1.
\]

Consider the signed mass-one group-algebra element
\[
  p=-2u^2+5u-2.
\]
Its total mass is \(p(1)=-2+5-2=1\), and
\[
  p(1)=1,\qquad p(2)=0,\qquad p(1/2)=0.
\]
Hence
\[
  \pi(p)=\operatorname{diag}(1,0,0),
\]
the projection onto \(E^\pi\). Also \(\pi(p)^2=\pi(p)\), so \(p\) is
idempotent in the quotient completion \(\cC_{\cF}(\Z)\). Thus \(p\) is a
Kazhdan projection in the sense of the source paper.

It remains to prove that no positive probability approximation exists. Let
\(\mu\) be a finitely supported probability measure on \(\Z\). Then
\[
  \pi(\mu)=\operatorname{diag}
  \left(1,\sum_n \mu(n)2^n,\sum_n \mu(n)2^{-n}\right).
\]
If a net \(\mu_i\) of such measures converged to \(p\) in
\(\cC_{\cF}(\Z)\), then both
\[
  A_i=\sum_n \mu_i(n)2^n,\qquad
  B_i=\sum_n \mu_i(n)2^{-n}
\]
would tend to zero. This is impossible, because for every \(i\),
\[
  A_i+B_i
  =\sum_n \mu_i(n)(2^n+2^{-n})
  \geq 2\sum_n\mu_i(n)
  =2.
\]
Therefore \(p\) is not in the closure of the nonnegative probability measures.
\end{proof}

\section*{Verification Notes}

The construction uses the infinite compactly generated group \(\Z\), so it
fits the source's standing locally compact compactly generated setting. The
representation is strongly continuous because \(G\) is discrete. The family
\(\cF\) has one representation, and the generator norms are finite. The
completion is finite-dimensional: it is the image of \(\C[\Z]\) under
\[
  a\mapsto(\widehat a(1),\widehat a(2),\widehat a(1/2)).
\]
Thus all convergence statements reduce to convergence of the two scalar
moments at \(2\) and \(1/2\).

\section*{Limitations}

The counterexample is literal for arbitrary Banach-space representation
families. It is not an example from strong property (T), robust property (T),
or isometric representation theory. It also does not refute Proposition 3.13
of the source, which assumes an isometric family stable by complex
conjugation.

\section*{Novelty Check}

The run indexes were searched for arXiv:1604.01616, the title, Kazhdan
projection, positive Kazhdan projection, nonnegative functions, Banach property
\((T)\), fixed point spectrum, Laplacian, cohomology, and expander keywords.
No exact packet for this source question was found. Nearby indexed results
concern BFGM higher-rank fixed point properties, fixed point spectra, and
expander/superexpander questions.

External phrase searches on June 28, 2026 used the phrases
\texttt{positive Kazhdan projections}, \texttt{nonnegative functions},
\texttt{Kazhdan projection}, \texttt{closure of the nonnegative functions},
and the exact source sentence. They did not locate a published answer. The
novelty confidence is therefore bounded by phrase search, but the example is
self-contained.

\section*{Human Review Recommendation}

Human review should check the interpretation of the phrase ``in general'' in
Section 3.5. If it is read in the full generality of the paper's arbitrary
Banach-space representation families, the counterexample is complete. If the
intended question was restricted to isometric, self-adjoint, or naturally
arising strong-property-\((T)\) families, this packet should be reclassified as
a scope caveat rather than a resolution of that narrower problem.

\begin{thebibliography}{9}

\bibitem{deLaSalle2016}
M. de la Salle,
\emph{A local characterization of Kazhdan projections and applications},
arXiv:1604.01616.

\end{thebibliography}

\end{document}
